\paragraph{Conjugate-coordinate class: the Casimir of \(\mathfrak{so}(4,2)\) and its combinations as proper-time regulators in the four-matrix class (conditional; bounded no-go and conditional route).}
Keep the Poisson limit, the leaf \(\mathbb{CP}^{1,2}\) of \(\mathfrak{so}(4,2)\), \(T^a=M^{a4}\), \(D^a=\{T^a,\cdot\}\), \(J^{AB}=\{M^{AB},\cdot\}\), \(\Box=\eta_{ab}D^aD^b\), the vector operator \(H_V\) with \(\mu^2=-3\) and the notation of the preceding paragraphs. Write \(C_{42}\), \(C_{41}\), \(C_{31}\) for \(\tfrac12\eta_{AC}\eta_{BD}J^{AB}J^{CD}\) summed over the indices \(0,\dots,5\), \(0,\dots,4\), \(0,\dots,3\), and \(\sigma(\cdot)\) for the quadratic form \(\sum w\,(\xi\!\cdot\!v)^2\) of the cotangent vector \(\xi\) of an operator \(\sum w\,v^2\) of vector fields \(v\) (the principal symbol is its negative).

\emph{Exact relations.} In the polynomial Poisson algebra, for both signs of the bracket, \(\Box=C_{41}-C_{31}\), \(\mu^2=-3\) and, with \(C_G\) the Casimir of the total \(\mathfrak{so}(3,1)\) action on vector-valued functions (orbital part plus the vector index, the sign of the index action correlated with the sign of the bracket) and \(C_\Sigma\) that of the vector index,
\[
H_V=C_G-C_{41}-C_\Sigma+\mu^2 ,
\]
the unique solution of the linear fit \(H_V=xC_G+yC_{31}+zC_{41}+wC_{42}+uC_\Sigma+\mu^2\) with \(w=0\). Hence \(C_{42}\) does not occur in \(H_V\), and \(C_{42}-C_{41}=-\sum_{A\le4}\eta_{AA}(J^{A5})^2\). Since \(C_{42}\) is central in \(U(\mathfrak{so}(4,2))\) and \(C_{41}\) commutes with \(\mathfrak{so}(4,1)\), both commute with \(H_V\), with \(\Box\) and with the ghost operator \(-\Box\); a function of \(C_{42}\) and \(C_{41}\) therefore commutes with every operator of the class and no ordering ambiguity arises.

\emph{Symbols.} At \(40\) random orbit points the form \(\sigma(C_{42})\) on the six-dimensional cotangent space has signature \((2,4)\), and so has \(\sigma(C_{41})\). The pencil \(\sigma(C_{41})v=\lambda\,\sigma(C_{42})v\) has the eigenvalues \(\tfrac12\) (four times) and \(1\) (twice) at \(30\) points (deviation \(3\cdot10^{-13}\)); on the four-dimensional eigenspace (the span of the \(dy^C\), \(y^C=M^{C5}\)) \(\sigma(C_{42})\) is negative definite and on the two-dimensional (vertical) eigenspace positive definite. The base symbol \(-2\ell^2g^{-1}\) of the preceding paragraph is the first of these; the fibre carries the opposite sign. Write \(\sigma(C_{42})=-\alpha q_b+\beta q_f\) in this splitting, so that \(\sigma(C_{41})=-\tfrac12\alpha q_b+\beta q_f\). The operators \(C_{41}-C_{42}\) (the five-matrix \(\Box\), purely horizontal) and \(2C_{41}-C_{42}\) (purely vertical) are the two degenerate members of the family
\[
\Delta_\kappa=\kappa\,C_{42}-C_{41},\qquad \sigma(\Delta_\kappa)=\alpha(\kappa-\tfrac12)q_b+\beta(1-\kappa)q_f ,
\]
which is positive definite exactly for \(\tfrac12<\kappa<1\) (checked at \(15\) points for nine values of \(\kappa\)); the value \(\kappa=\tfrac34\) gives \(\tfrac14(\alpha q_b+\beta q_f)\), the absolute value of \(\sigma(C_{42})\) in the splitting. The horizontal planes \(\operatorname{span}\{D^a\}\) of the four-matrix class and \(\operatorname{span}\{J^{A5}\}\) of the five-matrix class are different four-planes (their sum has dimension \(6\) at six points).

\emph{Consequences.} (i) No function of \(C_{42}\) alone is an elliptic regulator: the level sets of an indefinite form in the cotangent fibre are non-compact, so \(\int\mathrm d^6\xi\,\chi(\sigma(C_{42})/\LUV^2)\) diverges for every \(\chi\) of compact support, and a cutoff of \(C_{42}\) leaves the null cone \(q_b=\beta q_f/\alpha\) unregulated; this is a statement about the leading symbol. (ii) \(\Delta_\kappa\), \(\tfrac12<\kappa<1\), is built from the matrix data \(M^{AB}\) and the contraction \(\eta_{AB}\) only, commutes with \(H_V\), \(\Box\) and the ghost operator, and has a definite symbol on the whole leaf. For a cutoff function \(\chi\) of compact support, the density \((2\pi)^{-6}\!\int\mathrm d^6\xi\,\chi(\sigma(\Delta_\kappa)/\LUV^2)\,F(\sigma(H_V))\) is finite at every point for every continuous \(F\) (at the level of the leading symbol); in particular the central variable \(\zeta\) of the preceding paragraphs is confined to a ball \(|\zeta|\lesssim\LUV\), which regulates the \(\zeta\)-integral left divergent there. It is a different regulator from a function of \(\Box_X\), and it uses the generators \(J^{A5}\), which are not in \(U(\mathfrak{so}(4,1)')\); it is an additional element of the regularisation and not derived from the action. (iii) The volume of the window \(\{\sigma(\Delta_\kappa)\le\LUV^2\}\) is proportional to \(\LUV^6[(\kappa-\tfrac12)^2(1-\kappa)]^{-1}\) (checked against \(\det^{-1/2}\) at five values of \(\kappa\)), so the leading coefficient depends on \(\kappa\) and diverges at both ends of the interval, where the horizontal and the vertical directions respectively are unregulated; no value of \(\kappa\) is singled out by the algebra, and the quadratic coefficient, which is the one an Einstein--Hilbert term needs, is not independent of the choice of \(\kappa\) and \(\chi\) without a further argument. (iv) For the Lorentzian \(\eta\) the form \(\sigma(\Box)\) has signature \((3,1)\) with two null directions, and in the window of \(\Delta_{3/4}\) the generalised eigenvalues of \((\sigma(\Box),\sigma(\Delta_{3/4}))\) are \(-2\) (at eight points, the same value at each) and three positive values that grow with \((y^4)^2\) (so the negative bound is uniform in position and the positive range is not); the symbol of \(H_V\) is therefore bounded by \(\mp\) a multiple of \(\LUV^2\) on the window but is negative on part of it, so the symbol of \(e^{-sH_V}\) is bounded by \(e^{2s\LUV^2}\) there and the Euclidean proper-time integral at large \(s\) still diverges. The window removes the divergence at large \(\zeta\), not the negative direction: for the Euclidean contraction \(\eta\to\delta\) of the preceding paragraphs the lowest Zeeman level \(-|f_1-f_2|\) is bounded in the window of \(\Delta_{3/4}\) by \(d_*\LUV\) (in the normalisation of the forms above) with \(1.7\le d_*\le5.6\) at six points, linear in the cutoff.

\emph{Compact check.} For the compact form \(\mathfrak{so}(6)\) with \(\delta\), \(X^a=M^{a5}\), \(a=1,\dots,4\), \(m^2=-3\), on the \(N\)-th symmetric power (the matrices \(\operatorname{End}(\mathrm{Sym}^N\mathbb C^4)=\bigoplus_{k\le N}(k,k,0)\) of \(\mathfrak{so}(6)\), so that a cutoff on \(C_{\mathfrak{so}(6)}=2k^2+6k\) is the truncation of the matrix algebra), the gauge-fixed Hessian \(H+G^\dagger G\), \(Ga=\sum_a[X^a,a^a]\), has the spectrum
\[
\lambda=c_5(a,b)-c_4(c',d')+3+m^2,
\]
with \(c_5(a,b)=a(a+3)+b(b+1)\), \(c_4(c,d)=c(c+2)+d^2\), \(a=k\) the label of \(\mathfrak{so}(6)\), \(0\le b\le a\), \((c,d)\) the \(\mathfrak{so}(4)\) content of \((a,b)\) and \((c',d')\) that of \((c,d)\otimes(1,0)\); this was checked against the numerical spectrum for \(N=1,2,3\) (differences below \(10^{-12}\)). The lowest eigenvalue in the window \(k\le A\) is \(-(A+3)\), the number of negative modes is \(38,124,344,750\) for \(A=1,\dots,4\), and their fraction of the \(4\binom{A+3}{3}^2\) modes falls to \(2.5\cdot10^{-3}\) at \(A=48\) (local exponents \(3.9\) and \(5.7\), approaching \(4\) and \(6\); not proved). The cutoff by the Casimir therefore bounds the negative direction by a quantity growing linearly in the cutoff and removes none of it.

\emph{Eight matrices.} For \(T^a\) and \(y^a=M^{a5}\), \(a=0,\dots,3\), with the contraction \(\eta\) on \(T\) and \(\lambda\eta\) on \(y\), exact stationarity requires the bilinear coefficients \(\lambda-3\) for \(T\) and \(3\lambda-1\) for \(y\), equal only at \(\lambda=-1\) (common value \(-4\)). The contraction form on the tangent space has signature \((3,3)\) at \(\lambda=-1\) and \(-100\le\lambda\le-0.01\) (thirty points each), \((5,1)\) at \(\lambda=0.01,0.1,10\) and is degenerate of rank \(4\) at \(\lambda=1\); the horizontal span has rank \(6\) throughout, so there is no centre variable. A Lorentzian \((5,1)\) induced form with no centre thus occurs for \(\lambda>0\), \(\lambda\ne1\), only with two independent bilinear coefficients; the Hessian of that set is not examined.

This is a statement for the Poisson limit, one loop, the leading symbol, the four-matrix class with the bilinear term, the regulators that are functions of \(C_{42}\) and \(C_{41}\), and the sampled points; the signatures and the pencil are numerical, the relations in the first paragraph exact. It does not treat fermions, the Moyal terms, a Lorentzian \(i\epsilon\) integral, the identification of the window with any admissibility condition of the fibre, or the subleading symbol, it does not compute the weight-two coefficient or its variation under \(A\), and neither the sign nor the size of \(w\) follows. It identifies no induced metric with the emergent metric of Chapter~\ref{ch:vierbein-emergence} and, since a Euclidean coefficient is not a Lorentzian one, leaves Assumption~\ref{ass:lorentzian} in place (Problem~E). This changes no ledger label; Problem~F remains open.
